\documentclass[11pt]{article}
\usepackage[margin=1in]{geometry}
\usepackage{amsmath,amssymb,amsthm}
\usepackage[hidelinks]{hyperref}
\newtheorem{theorem}{Theorem}
\newcommand{\C}{\mathbb C}
\newcommand{\A}{\mathcal A}
\newcommand{\norm}[1]{\left\lVert#1\right\rVert}
\title{No finite C*-Smale mean-value constant at degree three}
\author{Research proof and Lean formalization}
\date{October 27, 2026}
\begin{document}
\maketitle
\begin{abstract}
Over the single commutative C*-algebra $C([0,1],\C)$, no finite constant
controls the Smale norm divided difference for all split monic cubics.
For every proposed constant, a continuous cubic family has an invertible
derivative at the test element and exactly two existing critical elements,
both violating the bound. This uniform impossibility statement is stronger
than a fixed counterexample to the constant-one conjecture.
\end{abstract}

\begin{theorem}\label{main}
Let $\A=C([0,1],\C)$ with its supremum norm. For every $C\geq0$ there are
$a_1,a_2,a_3,z\in\A$ such that the monic cubic
$P(w)=\prod_{j=1}^3(w-a_j)$ has $P'(z)$ invertible and exactly two distinct
critical elements. Every critical element $b\in\A$ satisfies
\begin{equation}\label{violation}
 \frac{\norm{P(z)-P(b)}}{\norm{z-b}}>C\norm{P'(z)}.
\end{equation}
Here $P'$ is the literal sum of products obtained by deleting one factor.
\end{theorem}

\begin{proof}
Write $y(t)=2t-1$. For $0<\varepsilon\leq1/16$ put
\[
 q_\varepsilon(t)=\sqrt{y(t)^2+\varepsilon},\qquad
 P_\varepsilon(w)=w(w-\sqrt3 q_\varepsilon)(w+\sqrt3 q_\varepsilon)
                 =w^3-3q_\varepsilon^2w,\qquad z=y.
\]
Then
\begin{equation}\label{derivative}
 P_\varepsilon'(w)=3(w-q_\varepsilon)(w+q_\varepsilon),\qquad
 P_\varepsilon'(y)=-3\varepsilon\,1.
\end{equation}
The latter is invertible. Both $q_\varepsilon$ and $-q_\varepsilon$ are
critical. Conversely, scalar evaluation of $P_\varepsilon'(b)=0$ shows
$b(t)=\pm q_\varepsilon(t)$ at each point. The quotient
$b/q_\varepsilon$ is continuous into $\{-1,1\}$, since $q_\varepsilon>0$.
Connectedness forces one global sign, so these are exactly the two distinct
critical elements. No cancellation in a function ring is being assumed.

For either critical element,
\begin{equation}\label{uniform}
 0<\norm{y-b}\leq\norm y+\norm{q_\varepsilon}\leq\frac94,
 \qquad \norm{P_\varepsilon'(y)}=3\varepsilon.
\end{equation}
The denominator is nonzero because $y=b$ would contradict
\eqref{derivative}. With $r=\sqrt{1+\varepsilon}$, evaluation at the
opposite endpoints gives
\[
 (P_\varepsilon(y)-P_\varepsilon(q_\varepsilon))(0)
       =2+3\varepsilon+2r^3>2,
\]
\[
 (P_\varepsilon(y)-P_\varepsilon(-q_\varepsilon))(1)
       =-(2+3\varepsilon+2r^3).
\]
Consequently $\norm{P_\varepsilon(y)-P_\varepsilon(b)}>2$ for every
critical element, uniformly in $\varepsilon$.

Given $C\geq0$, choose
\begin{equation}\label{choice}
 \varepsilon=\frac1{16(C+1)}.
\end{equation}
This is strictly positive and at most $1/16$, and
\[
 C\norm{P_\varepsilon'(y)}\norm{y-b}
 \leq\frac{27}{4}C\varepsilon
 =\frac{27C}{64(C+1)}\leq\frac{27}{64}<2.
\]
Together with the value gap this proves \eqref{violation}.
\end{proof}

\paragraph{Relation to the source and prior refutations.}
Krishna \cite[Theorems~1.1--1.2]{Krishna} recalls the scalar Smale theorem
with universal constant $4$, and Conjecture~2.1 proposes a C*-algebraic
constant-one extension. Public scalar counterexamples, including
\cite[Theorem~1.3]{JatarNg}, already refute that universal conjecture because
$\C$ is itself a commutative C*-algebra. We do not claim its first resolution.
Theorem~\ref{main} rules out every finite constant at the fixed degree three
over one fixed function algebra; the scalar constant $4$ remains valid.
The source's degree-two theorem is known and receives no credit here.

\paragraph{Formal scope and review status.}
The Lean certificate has 49 declarations in five modules, including four
reused continuous-selection lemmas. The closed endpoint
\[
\texttt{CStarSmaleUnbounded.infinite\_cubic\_constant}
\]
contains the universal quantifier over $C$, the actual monic degree-three
product, invertibility of the test derivative, exact two-element critical
set and strict norm violation at every critical element. The second endpoint
\texttt{no\_finite\_cubic\_bound} negates the existence of a nonnegative
universal cubic bound. All source product/derivative and norm bridges,
\eqref{choice} and its positivity are proved. There is no residual small-
parameter or selection hypothesis.

The isolated strict replay passes under Lean~4.34.1 and mathlib revision
\[\texttt{d13f23b723b8a846827a245b89c10fc7d3f11612},\]
using only \texttt{propext}, \texttt{Classical.choice} and \texttt{Quot.sound}.
Independent final source, mathematical, formal, manuscript and bounded
novelty review passes. Eligibility is solely as a new infinite-class theorem,
not credit for an unresolved source question. Acceptance also requires a fresh
replay of the extracted proof archive. The rejected fixed constant-one witness
is not counted; the parameter family is one combined outcome.
No worldwide priority claim is made.

\begin{thebibliography}{2}
\bibitem{Krishna} K.~Mahesh Krishna,
\emph{C*-algebraic Smale mean value conjecture and Dubinin--Sugawa dual mean
value conjecture}, Vestnik KRAUNC. Fiz.-Mat. Nauki \textbf{54}(2) (2026), 38--47.
\href{https://doi.org/10.26117/2079-6641-2026-54-2-38-47}{doi:10.26117/2079-6641-2026-54-2-38-47}.
\bibitem{JatarNg} A.~Jatar and T.-W.~Ng,
\emph{Smale's mean value constant and admissible univalent Blaschke pairs},
\href{https://arxiv.org/abs/2610.21472v2}{arXiv:2610.21472v2} (2026).
\end{thebibliography}
\end{document}
